%%%%%%%%%%%%%%%%%%%%%%%%%%%%%%%%%%%%%%%%%%%%%%%%%%%%%%%%%%%%%%%%%%%
% ORGANIGRAMMES
%%%%%%%%%%%%%%%%%%%%%%%%%%%%%%%%%%%%%%%%%%%%%%%%%%%%%%%%%%%%%%%%%%%

% boxes
\newcommand{\boxvalue}[3]% #1 : abscisse, #2 ordonnée, #3 comment
{
    \draw[rounded corners = 4pt,thick,fill = cyan, color = cyan,fill opacity= 0.2] (#1-1,#2-0.5) rectangle (#1+1,#2+0.5)
    node[text = black,text opacity = 1,midway]{#3};
}

\newcommand{\boxvaluelarge}[3]% #1 : abscisse, #2 ordonnée, #3 comment
{
    \draw[rounded corners = 4pt,thick,fill = cyan, color = cyan,fill opacity= 0.2] (#1-2,#2-0.5) rectangle (#1+2,#2+0.5)
    node[text = black,text opacity = 1,midway]{#3};
}

\newcommand{\boxaction}[4]% #1 : abscisse, #2 ordonnée, #3 comment, #4 color
{
    \draw[rounded corners = 4pt,thick, color = #4] (#1-1,#2-0.5) rectangle (#1+1,#2+0.5)
    node[midway]{#3};
}

\newcommand{\boxactionlarge}[3]% #1 : abscisse, #2 ordonnée, #3 comment
{
    \draw[rounded corners = 4pt,thick] (#1-2,#2-0.5) rectangle (#1+2,#2+0.5)
    node[midway]{#3};
}

\newcommand{\boxactionthick}[3]% #1 : abscisse, #2 ordonnée, #3 comment
{
    \draw[rounded corners = 4pt,thick] (#1-2,#2-1) rectangle (#1+2,#2+1)
    node[midway]{#3};
}

\newcommand{\boxquestion}[3]% #1 : abscisse, #2 ordonnée, #3 comment
{
    \draw[rounded corners = 4pt,thick,fill = green, color = green,fill opacity= 0.2] (#1-1,#2-0.5) rectangle (#1+1,#2+0.5)
    node[text = black,text opacity = 1,midway]{#3};
}

\newcommand{\boxquestionlarge}[3]% #1 : abscisse, #2 ordonnée, #3 comment
{
    \draw[rounded corners = 4pt,thick,fill = green, color = green,fill opacity= 0.2] (#1-2,#2-0.5) rectangle (#1+2,#2+0.5)
    node[text = black,text opacity = 1,midway]{#3};
}

\newcommand{\boxbreak}[5]% #1 : abscisse center, #2 ordonnée center, #3 demi petit axe, #4 demi grand axe, #5 comment
{
    \draw[rounded corners = 4pt,thick,fill = red, color = red,fill opacity= 0.35] (#1-2,#2-0.5) ellipse  (#3 cm and #4 cm)
    node[text = black,text opacity = 1]{#5};
}

\newcommand{\comparator}[4]% #1 : abscisse, #2 ordonnée, #3 radius
{
    \draw[thick] (#1,#2) circle (#3) node{\textcolor{#4}{$+$}};
}

% links
\newcommand{\linkright}[5]% #1 : abscisse, #2 ordonnée, #3 length, #4 comment, #5 color
{
    \draw[thick,->] (#1,#2) --++ (#3,0) node[midway,above]{#4};
}

\newcommand{\linkdown}[5]% #1 : abscisse, #2 ordonnée, #3 length, #4 comment
{
    \draw[thick,->, color = #5] (#1,#2) --++ (0,-#3) node[midway,right]{#4};
}


\newcommand{\tmpprofile}[2]% #1 : abscisse, #2 ordonnée
{
	\draw[dashed] (#1,#2) ellipse (1.75cm and 1.45cm);
	\draw[thick,->] (#1-0.9,#2-0.6) --++ (2,0) node[midway,below = 5pt]{\footnotesize distance};
	\draw[thick,->] (#1-0.9,#2-0.6) --++ (0,1.5) node[rotate = 90,midway,above = 5pt]{\footnotesize power};
	\draw[thick] (#1-0.75,#2+0.4) --++ (0.625,-0.75) --++ (0,0.75) --++ (0.625,-0.75) --++ (0,0.75) --++ (0.3125,-0.375);
}